\documentclass[12pt]{article}
\usepackage[utf8]{inputenc}
\usepackage{amsmath, amssymb}
\usepackage{geometry}
\geometry{margin=1in}

\title{Solving for $L_M$ Step-by-Step}
\author{Specific-Factors Model}
\date{}

\begin{document}
\maketitle

\section*{Setup}

Home produces Manufactures ($M$) and Food ($F$).  
Labor is mobile between sectors, while capital is specific to $M$ and land to $F$.  

\[
Q_M = K^{1/2}L_M^{1/2}, 
\qquad 
Q_F = T^{1/2}L_F^{1/2},
\qquad 
L_M + L_F = L.
\]

Endowments:
\[
K=100, \quad T=64, \quad L=100.
\]

Prices: initially $P_M=2$, $P_F=1$.  
We will later examine the case when $P_M=3$.

\bigskip

\noindent The equilibrium condition (labor mobility):
\[
P_M \cdot MPL_M = P_F \cdot MPL_F,
\]
where
\[
MPL_M = \frac{\partial Q_M}{\partial L_M} = \frac{1}{2}K^{1/2}L_M^{-1/2}, 
\qquad
MPL_F = \frac{\partial Q_F}{\partial L_F} = \frac{1}{2}T^{1/2}L_F^{-1/2}.
\]

Since $K^{1/2}=10$ and $T^{1/2}=8$:
\[
MPL_M = 5L_M^{-1/2}, \qquad MPL_F = 4L_F^{-1/2}.
\]

Also $L_F = 100 - L_M$.

---

\section*{Case 1: $P_M = 2$, $P_F = 1$}

\noindent Step 1: Wage equalization condition
\[
2 \cdot 5L_M^{-1/2} = 1 \cdot 4(100 - L_M)^{-1/2}.
\]
Simplify:
\[
10L_M^{-1/2} = 4(100 - L_M)^{-1/2}.
\]

\bigskip
\noindent Step 2: Square both sides to remove the square roots
\[
100L_M^{-1} = 16(100 - L_M)^{-1}.
\]

\bigskip
\noindent Step 3: Cross-multiply
\[
100(100 - L_M) = 16L_M.
\]

\bigskip
\noindent Step 4: Expand and solve for $L_M$
\[
10{,}000 - 100L_M = 16L_M \quad \Rightarrow \quad 10{,}000 = 116L_M.
\]
\[
\boxed{L_M = 86.21, \quad L_F = 13.79.}
\]

\bigskip
\noindent Step 5: Compute the common wage
\[
w = P_M \cdot MPL_M = 2 \cdot 5L_M^{-1/2} = 10L_M^{-1/2}.
\]
\[
w = 10 / \sqrt{86.21} \approx 10 / 9.29 \approx \boxed{1.08.}
\]

---

\section*{Case 2: $P_M = 3$, $P_F = 1$}

\noindent Step 1: Wage equalization
\[
3 \cdot 5L_M^{-1/2} = 1 \cdot 4(100 - L_M)^{-1/2},
\]
\[
15L_M^{-1/2} = 4(100 - L_M)^{-1/2}.
\]

\bigskip
\noindent Step 2: Square both sides
\[
225L_M^{-1} = 16(100 - L_M)^{-1}.
\]

\bigskip
\noindent Step 3: Cross-multiply
\[
225(100 - L_M) = 16L_M.
\]
\[
22{,}500 - 225L_M = 16L_M.
\]

\bigskip
\noindent Step 4: Solve for $L_M$
\[
22{,}500 = 241L_M \quad \Rightarrow \quad \boxed{L_M = 93.36, \quad L_F = 6.64.}
\]

\bigskip
\noindent Step 5: Compute new wage
\[
w = P_M \cdot MPL_M = 3 \cdot 5L_M^{-1/2} = 15L_M^{-1/2}.
\]
\[
w = 15 / \sqrt{93.36} = 15 / 9.66 \approx \boxed{1.55.}
\]

---

\section*{Interpretation}

\begin{itemize}
\item As $P_M$ rises, $L_M$ increases (from 86.2 to 93.4), meaning more labor moves to manufacturing.  
\item The wage rises (from 1.08 to 1.55) but less than the increase in $P_M$ (which rose 50\%).  
\item Capital (specific to $M$) gains the most because it benefits from both the price increase and more labor.  
\item Land (specific to $F$) loses because $L_F$ falls and $P_F$ is unchanged.  
\item Workers have an ambiguous outcome: their real wage in terms of $F$ rises, but in terms of $M$ falls.
\end{itemize}

\end{document}
